
\section{本章内容}

本章回答"如何系统地设计一个本体"，介绍：
\begin{itemize}
\item 能力问题（Competency Questions, CQ）驱动的需求定义
\item Ontology Development 101 七步法
\item METHONTOLOGY 生命周期方法
\item NeOn 场景化方法概览
\item OntoClean 本体质量评估（元属性标注）
\end{itemize}

\section{文件说明}

\begin{center}
\begin{tabularx}{\linewidth}{@{}XX@{}}
\toprule
\textbf{文件} & \textbf{内容} \\
\midrule
\texttt{competency-questions.txt} & 制造领域能力问题设计：CQ\(\rightarrow\)本体元素\(\rightarrow\)验证查询 \\
\texttt{ontology-101-process.txt} & Ontology 101 七步法完整走查（制造本体） \\
\texttt{methontology-lifecycle.txt} & METHONTOLOGY 生命周期与概念化产物 \\
\texttt{ontoclean-evaluation.txt} & OntoClean 元属性标注与建模错误检查 \\
\bottomrule
\end{tabularx}
\end{center}

\section{方法论对比}

\begin{center}
\begin{tabularx}{\linewidth}{@{}llXX@{}}
\toprule
\textbf{方法论} & \textbf{提出时间} & \textbf{特点} & \textbf{适用场景} \\
\midrule
Ontology 101 & 2001 & 七步迭代，轻量易上手 & 教学、中小型本体 \\
METHONTOLOGY & 1997 & 完整生命周期+产物模板 & 企业级、需文档化交付 \\
NeOn & 2009 & 九种场景，强调复用 & 网络化、多本体集成 \\
OntoClean & 2002 & 元属性形式化评估 & 类层次质量审查 \\
\bottomrule
\end{tabularx}
\end{center}

\section{关键概念}

\begin{itemize}
\item \textbf{能力问题 (CQ)}: 本体建成后必须能回答的问题，是范围定义与验收标准
\item \textbf{概念化 (Conceptualization)}: 从术语表到概念分类树、关系表的中间产物
\item \textbf{刚性 (Rigidity)}: OntoClean 元属性，区分本质类型与角色/状态
\item \textbf{本体复用 (Reuse)}: 优先采用顶层本体（BFO/DOLCE）与领域本体（IOF），见附录A
\end{itemize}
